\section{Conclusiones y Trabajos Futuros}

\subsection{Conclusiones}

Tras un exhaustivo proceso de diseño, implementación y evaluación, emerge una conclusión contundente: la compresión geoespacial basada en inteligencia artificial constituye una herramienta madura y altamente efectiva para enfrentar las restricciones de ancho de banda en misiones de exploración remota. El marco propuesto logró reducciones de bitrate superiores al 74\% respecto al estándar VTM, manteniendo al mismo tiempo una degradación mínima en tareas downstream críticas —por debajo del 5\% en la mayoría de escenarios evaluados—.

Molliere et al. \cite{Molliere2025_Learned_EO_Compression} subrayan en su revisión la necesidad de enfoques que equilibren compresión agresiva con preservación de valor analítico en Earth Observation. Nuestro trabajo responde directamente a esa necesidad mediante la integración de características multiespectrales y mecanismos semánticos, logrando un balance práctico que supera a varios métodos del estado del arte.

Particularmente gratificante resulta comprobar que las mejoras no se limitan al dominio numérico. Las reconstrucciones preservan detalles estructurales y texturales esenciales para clasificación de cobertura, detección de anomalías y análisis de recursos, validando así la hipótesis central del proyecto.

Aunque los resultados son prometedores, cabe matizar que ningún sistema actual elimina por completo el trade-off inherente entre tasa de compresión y fidelidad. Esta consciencia crítica refuerza la importancia de evaluar siempre la utilidad final de los datos en el contexto operativo específico.

\subsection{Trabajos Futuros}

Lejos de ser un asunto resuelto, el campo sigue ofreciendo oportunidades significativas de avance. Ye et al. \cite{Ye2025_MAPC} demuestran el potencial de enfoques asistidos por información auxiliar para alcanzar bitrates extremadamente bajos. Una extensión natural de nuestro trabajo consiste en incorporar mecanismos map-assisted o semantic-guided que permitan compresión adaptativa según la misión y la región observada.

Entre las líneas prioritarias destacan:

\begin{itemize}
    \item Implementación onboard de versiones ligeras del modelo, optimizadas para el entorno de computación embarcada en satélites con restricciones severas de potencia y memoria.
    \item Exploración de modelos de difusión y arquitecturas generativas para mejorar la reconstrucción de texturas complejas en rangos de bitrate muy bajos.
    \item Desarrollo de sistemas de compresión continua y adaptativa que aprovechen información temporal entre pasadas sucesivas.
    \item Integración multimodal que combine datos ópticos, SAR y otras fuentes para mayor robustez en condiciones adversas.
    \item Estudios de viabilidad económica y operativa en misiones reales, incluyendo el impacto en costos de downlink y tiempos de respuesta.
\end{itemize}

Molliere et al. \cite{Molliere2025_Learned_EO_Compression} sugieren además la necesidad de avanzar hacia estándares abiertos de compresión aprendida para Earth Observation. Creemos que contribuir a esa estandarización debe formar parte de los esfuerzos futuros de la comunidad.

En definitiva, el trabajo presentado marca un avance concreto hacia una transmisión más eficiente de datos geoespaciales. Queda aún un camino interesante por recorrer. La combinación de rigor técnico con una visión orientada a aplicaciones reales ofrece, en nuestra experiencia, el mejor camino para transformar la compresión IA de un tema académico en un componente estándar de las futuras misiones de exploración espacial y monitoreo terrestre.